\section{All coefficient test rings for the original framed points}

\paragraph{The assertion proved.}
Fix the original \(E,H,D,l,X_o,c\), and put \(K=c.F\) and
\(X=X_o\times_{\operatorname{Spec}K_A(l)}\operatorname{Spec}K\).
The supplied \texttt{Moduli} record, including its universal object,
smoothness and quasiprojectivity, is retained.  The only additional
section-theoretic premise is the original
\(\mathrm{BC}=\texttt{Out.SecBaseChangeSt}\).  This premise is applied
to a locally Noetherian base such as \(X_o\) or \(X\).  Its conclusion
allows \emph{every} commutative coefficient ring on the other side of
the base change.  We never replace it by a flat-base-change assertion,
and we do not assume the stronger predicate
\texttt{SecBaseChangeAnySt} with arbitrary original base schemes.

The two principal conclusions of this chapter are
\begin{align}
 &\text{for every commutative }K\text{-algebra }S
   \text{ and }y\in\operatorname{FrPt}(S),\notag\\
 &\qquad\exists\ S_0\subseteq S\text{ a finitely generated }K
       \text{-subalgebra},\notag\\
 &\qquad y_0\in\operatorname{FrPt}(S_0),\quad
       \operatorname{FrMap}(S_0\hookrightarrow S,y_0,y);       \tag{TDR.1}\\
 &\text{if }S\text{ is reduced, there are }n,I,\phi,y'\text{ with}\notag\\
 &\qquad I\subseteq K[t_1,\ldots,t_n]\text{ radical},\quad
       \phi:K[t_1,\ldots,t_n]/I\longrightarrow S,\notag\\
 &\qquad y'\in\operatorname{FrPt}(K[t_1,\ldots,t_n]/I),\notag\\
 &\qquad
       \operatorname{FrMap}(\phi,y',y).                       \tag{TDR.2}
\end{align}
Thus (TDR.1) includes nonreduced, non-Noetherian and non-finite-type
\(S\), and (TDR.2) is the literal \texttt{AllTest} conclusion.
The algebra map \(\phi\), tangent vectors and section are part of the
conclusion; the result is not merely equality of isomorphism classes
of abelian varieties or a statement about field-valued points.

\paragraph{Geometric dictionary.}
Let \(\pi:A\to X\) and \(L\) be the pullbacks of the actual universal
abelian scheme and its prescribed principal theta line.  Write
\(m_\tau=D.\mathrm{mult}(\tau)\), for
\(\tau:E\hookrightarrow\mathbb C\), and
\(\tau_K:E\to K\) for the same embedding with codomain restricted
using \(c.\mathrm{conj}\).  The finite set of all these embeddings,
including those with \(m_\tau=0\), is retained.  A framed point has an
underlying map \(x:\operatorname{Spec}S\to X\), ordered tangent
eigenbases \((u_{\tau,i})_{i<m_\tau}\), and a section \(s\) which is
an \(S\)-module generator of \(H^0(A_x,L_x)\).  The source notation
\(\operatorname{PPx}(x)\) is the same line, expressed as its
one-fold \texttt{LPow} on the canonical pullback.  We identify these
lines only by that canonical pullback isomorphism.

The standard algebraic foundations used below are the description of
tangent vectors by dual numbers, Lie base change for a smooth group
scheme, the sheaf properties of regular functions and line-bundle
sections, and affine scheme/ring duality.  Frame schemes and all
finite-descent arguments needed here are constructed below.  No
analytic uniformization, choice of Haar measure, or new result about
automorphic forms is used.

\subsection{The actual frame functor over every coefficient algebra}

% BEGIN TDR1
\paragraph{TDR1. Tangent eigenbases are frames of fixed vector bundles.}
\label{claim:TDR1}
There are finite locally free bundles \(V_\tau\) of rank \(m_\tau\)
on \(X\) such that, for every commutative \(K\)-algebra \(S\) and
\(x:\operatorname{Spec}S\to X\), the source's
\(\tau\)-eigen tangent vectors are exactly
\(\Gamma(\operatorname{Spec}S,x^*V_\tau)\).
Under this identification \texttt{IsCombAt} is ordinary linear
combination.  The source's \texttt{basis} condition is precisely an
isomorphism \(S^{m_\tau}\simeq\Gamma(\operatorname{Spec}S,x^*V_\tau)\).
These identifications commute with arbitrary \(K\)-algebra maps.

\begin{proof}
A morphism from \(\operatorname{Spec}S[\epsilon]/(\epsilon^2)\)
to a smooth group scheme which reduces to its identity is a section
of its relative Lie bundle.  On an affine chart, its pullback on
functions is \(a\mapsto e^*a+\epsilon\delta(a)\), where
\(\delta\) is an \(S\)-linear derivation at the identity.  Multiplying
two such group-valued points adds their derivations: terms involving
two infinitesimal factors vanish because \(\epsilon^2=0\).
Precomposition by \(\epsilon\mapsto b\epsilon\) multiplies the
derivation by \(b\).  Consequently the finite product in
\texttt{IsCombAt} is exactly the sum of the indicated scalar
multiples.  This also proves that the source \texttt{tan} predicate,
including both its structural and identity-section equations, is
the relative tangent condition just described.

The actual \(\mathcal O_E\)-action differentiates to the Lie bundle
\(V=\operatorname{Lie}(A/X)\).  It extends to the split algebra
\[
 E\otimes_{\mathbb Q}K\simeq\prod_{\tau:E\hookrightarrow\mathbb C}K,
 \qquad a\otimes b\longmapsto(\tau_K(a)b)_\tau .             \tag{TDR.3}
\]
Here \(K\) contains every conjugate of \(E\), and characteristic
zero makes \(E/\mathbb Q\) separable.  Explicitly, choose a primitive
element of \(E\); the Lagrange interpolation polynomials in its
distinct \(K\)-valued roots give mutually orthogonal idempotents
\(e_\tau\) with sum one.  All denominators are nonzero elements
of \(K\), hence remain units in every nonzero \(K\)-algebra.
The zero ring causes no exception: all the modules and assertions
over its empty spectrum are trivial.

Set \(V_\tau=e_\tau V\).  These are direct summands of a finite
locally free bundle, hence finite locally free.  The eigencondition
for every \(a\in\mathcal O_E\) is equivalent to membership in this
summand.  Indeed it implies the condition for \(E\), by clearing
integer denominators, and the interpolation idempotents then isolate
the required summand.  Conversely the idempotent summand has precisely
the required character.  The original determinant/multiplicity
condition in the universal object, evaluated over algebraic closures
of residue fields, says that its rank is \(m_\tau\).  A finite
locally free module has locally constant rank detected on these
fibers, so the same rank holds on \(X\).

Lie formation for a smooth group scheme commutes with every base
change: equivalently, pull back its locally free invariant
differentials and dualize.  Each \(e_\tau\) is an idempotent over
\(K\), so its summand also commutes with base change, including
nonflat maps.  Finally, unique expression of every eigenvector in
the ordered \(u_{\tau,i}\) is exactly bijectivity of the indicated
linear map from \(S^{m_\tau}\).  No argument using only generic
fibers is involved.  For \(m_\tau=0\), both modules are zero and
the same statement holds with the empty basis.
\end{proof}
% END TDR1

% BEGIN TDR2
\paragraph{TDR2. The theta-section factor is one invertible sheaf with arbitrary base change.}
\label{claim:TDR2}
Under \(\mathrm{BC}\), there is an invertible sheaf \(\Theta\) on
\(X\), canonically the direct-image sheaf of the actual \(L\), and
for every \(x:\operatorname{Spec}S\to X\) a natural isomorphism
\[
 \Gamma(\operatorname{Spec}S,x^*\Theta)
       \simeq H^0(A_x,L_x).                                 \tag{TDR.4}
\]
An element on the right satisfies the literal source \texttt{gen}
condition if and only if its element on the left trivializes
\(x^*\Theta\).  Such an element remains a generator after every
algebra map \(S\to T\), with no reducedness or flatness assumption.

\begin{proof}
The original moduli scheme is smooth over the field \(K_A(l)\),
so it is locally Noetherian; the same is true after the finite field
extension to \(K\).  Its universal object is the supplied principally
polarized abelian scheme.  Apply \(\mathrm{BC}\) with \(n=1\).
The index set of sections has size \(1^g=1\).  On a neighborhood
of each point it therefore supplies one section whose pullback along
\emph{every} map from an affine scheme is a basis of the pulled-back
section module.  Ordinary pullback of the section supplies the
\texttt{IsPullSec} witness in that predicate.  Its two asserted
conditions, linear independence and spanning, say exactly that the
section module is freely generated by this section.

On affine subopens these sections identify \(\pi_*L\) with the
structure sheaf.  On overlaps two such identifications differ by a
unique unit: each generator is a unique scalar multiple of the other,
and composing the two expressions gives the scalar product one.
The overlap units satisfy the cocycle law by the same uniqueness.
They define an invertible sheaf \(\Theta=\pi_*L\), rather than
merely a collection of one-dimensional vector spaces on closed points.

The same local generators and overlap units pull back along any
\(x\).  The defining assertion of \(\mathrm{BC}\) identifies the
pulled-back generator with a basis of \(H^0(A_x,L_x)\) on every
affine subopen mapping to a trivializing neighborhood.  Both sides
of (TDR.4) are sheaves for the Zariski topology on the base: sections
of the pulled-back line on \(A_x\) glue over preimages of base
opens.  The local identifications consequently glue to (TDR.4),
naturally in \(x\), and they are the canonical section pullback
maps.  This also proves the assertion when \(S\) is non-Noetherian.

On an affine scheme an invertible sheaf corresponds to a projective
module of rank one.  The source \texttt{gen} says that
\(a\mapsto a s\) is a bijection from \(S\) to that module,
with scalar multiplication written chartwise as \texttt{scal}.
This is exactly its trivialization.  Tensoring that isomorphism with
an arbitrary \(S\)-algebra \(T\) preserves the isomorphism;
no flatness is necessary.  The section obtained is the actual
pullback section by the preceding naturality.
\end{proof}
% END TDR2

% BEGIN TDR3
\paragraph{TDR3. Exact base change of the original framed-point record.}
\label{claim:TDR3}
For every \(K\)-algebra map \(\phi:S\to T\) and original
\(y\in\operatorname{FrPt}(S)\), there is a unique
\(\phi_*y\in\operatorname{FrPt}(T)\) with
\(\operatorname{FrMap}(\phi,y,\phi_*y)\).
It is functorial in \(\phi\).  Thus this relation is exactly
pullback of the underlying point, each ordered eigenframe and the
chosen section, with no scalar ambiguity.

\begin{proof}
Set \(x_T=x\circ\operatorname{Spec}\phi\), pull back each
basis of \(x^*V_\tau\), and pull back the generator of
\(x^*\Theta\).  TDR1 and TDR2 give all the tangent, eigenbasis
and section-generator fields of the new record.  In dual-number
notation its tangent vectors are literally
\[
 u'_{\tau,i}
   =u_{\tau,i}\circ\operatorname{Spec}(\phi[\epsilon]),
 \qquad
 \phi[\epsilon](a+b\epsilon)=\phi(a)+\phi(b)\epsilon .       \tag{TDR.5}
\]
For the section let \(p_m:A_{x_T}\to A_x\) be the canonical
pullback map.  It commutes with both projections in the definition
of \texttt{FrMap}.  On every original line-bundle chart the new
section is \(p_m^*s\), which is exactly the universally quantified
\texttt{appLE} equation in that definition.  If a chart is further
restricted, the equation remains true because restriction commutes
with pullback.  This verifies the literal relation, including its
last field, rather than only an isomorphism of section lines.

For uniqueness, the first two fields of \texttt{FrMap} determine
the point and all tangent vectors.  The two projection equations
uniquely determine \(p_m\) by the fiber-product property, and the
section equations determine every chart component of the section.
All remaining fields are properties of these data.  The identity
and composition laws now follow either directly from (TDR.5) and
section pullback or from this uniqueness.  In particular if
\(z=\phi_*y\), then \((\psi\phi)_*y=w\) if and only if
\(\psi_*z=w\).  This is the cancellation law used in the source
descent proof.
\end{proof}
% END TDR3

% BEGIN TDR4
\paragraph{TDR4. Change of frame, and descent of equality along injections.}
\label{claim:TDR4}
For two original framed points \(y,y'\) over the same underlying
\(x:\operatorname{Spec}S\to X\), there is a unique collection
\[
 g_\tau\in\operatorname{GL}_{m_\tau}(S),\qquad a\in S^\times
 \quad\text{such that}\quad
 u'_{\tau,i}=\sum_j(g_\tau)_{ij}u_{\tau,j},\qquad s'=a s.
                                                               \tag{TDR.6}
\]
These are the original \texttt{FrChg} equations.  They commute
with every coefficient algebra map.  If \(S\hookrightarrow T\)
is injective and \(y,y'\) have identical pullbacks over \(T\),
then \(y=y'\).

\begin{proof}
Each ordered basis of a free module gives unique coordinates for
the other.  The coordinate maps in both directions are inverse,
so the matrices in (TDR.6) have unit determinant.  TDR1 identifies
these coordinate equations with \texttt{IsCombAt}; the convention
is rows, exactly as in \texttt{FrChg}.  By TDR2, \(s\) and \(s'\)
are bases of the same rank-one module.  Expressing each in the other
gives unique scalars whose product is one.  The section equation
\(s'=a s\) is the source \texttt{SecBC} for the identity map.
This proves existence and uniqueness, including zero-rank
eigenspaces and arbitrary coefficient rings.

Applying a ring map to (TDR.6) gives the coordinates of the
pulled-back frames, by TDR3.  If those frames coincide, the images
of all matrices and of the scalar are respectively the identity
matrices and one.  Injectivity of the ring map forces the same
equations over \(S\).  Equation (TDR.6) then gives equality of
the original points.  The hypothesis that their underlying maps
\(x\) already agree is retained; we are not using an unjustified
assertion about arbitrary maps of schemes after an injective ring
extension.
\end{proof}
% END TDR4

% BEGIN TDR5
\paragraph{TDR5. Reduced affine charts with actual reference frames.}
\label{claim:TDR5}
Every point of \(X\) has an affine open neighborhood
\(j:\operatorname{Spec}R\hookrightarrow X\), where \(R\) is
a reduced finitely generated \(K\)-algebra, together with an
original framed point \(y_0\in\operatorname{FrPt}(R)\) whose
underlying map to \(X\) is \(j\).

\begin{proof}
The finitely many \(V_\tau\) and the line \(\Theta\) are locally
free, by TDR1 and TDR2.  Intersect trivializing neighborhoods for
all of them and then take an affine open neighborhood of the chosen
point.  The base \(X\) is smooth over \(K\), hence reduced and
locally of finite type.  Its affine open neighborhood is reduced;
its structural map is of finite type because it is affine and
locally of finite type.  Thus \(R\) has the asserted properties.

Choose the ordered bases of each \(V_\tau\) and a generator of
\(\Theta\) on this affine open.  TDR1 translates the bases into
the required dual-number tangent maps with the \(\tau_K\)-action,
and TDR2 translates the line generator into the actual section of
\(\operatorname{PPx}\), with its unique coefficient rule.
They give every field of \(y_0\).  Its underlying point is the
composition of \(j\) with the projection \(X\to X_o\), so the
two projection equations in \texttt{OpenChart} hold literally.
TDR12 below presents \(R\) as the original \texttt{TRing};
transport along that algebra isomorphism and TDR3 then give
\texttt{exists\_openChart\_testRing} with the same image in \(X\).
The finite-generation part of this proof does not assume that
the entire moduli scheme is covered by one affine chart.
\end{proof}
% END TDR5

% BEGIN TDR6
\paragraph{TDR6. The explicit universal frame algebra and its exact classifier.}
\label{claim:TDR6}
For a fixed affine chart and reference frame from TDR5, set
\[
 B=R[\,X_{\tau,ij},s\,]
       \left[\frac{1}{s\prod_\tau\det(X_\tau)}\right].       \tag{TDR.7}
\]
It is a reduced finitely generated \(K\)-algebra, and hence is an
original test ring up to a specified algebra isomorphism.  With
\(\iota:R\to B\), the universal matrices and scalar satisfy
\texttt{FrameRing}.  There is an actual universal frame
\(u\in\operatorname{FrPt}(B)\) such that, for every commutative
\(K\)-algebra \(T\), map \(\psi:R\to T\), and
\(z\in\operatorname{FrPt}(T)\) over \(j\circ\operatorname{Spec}\psi\),
there is a unique \(b:B\to T\) with
\[
 b\iota=\psi,\qquad\operatorname{FrMap}(b,u,z).             \tag{TDR.8}
\]

\begin{proof}
There are finitely many variables, including one variable for the
theta-section scalar.  Localization in (TDR.7) makes \(s\) and
each determinant units: each divides their product, and a factor
of a unit in a commutative ring is a unit.  Conversely, any choice
of invertible matrices and a unit in \(T\) sends this product to
a unit.  The polynomial universal property followed by the
localization universal property gives a unique algebra map
\(B\to T\) with these values and restriction \(\psi\).
This is the entire \texttt{FrameRing} statement, with its
universal quantifier over all \(T\).

A polynomial ring over a reduced ring is reduced.  One direct
verification is to embed the reduced ring into the product of
its integral-domain quotients by prime ideals and use the
coefficientwise embedding of polynomial rings; the target has no
nonzero nilpotents.  Localization also preserves reducedness:
if a localized element is nilpotent, multiplying the nilpotence
relation by a suitable denominator gives a nilpotent element in
the original ring, whence the localized element is zero.
Localization at one element is a finite-type operation, realized
by adjoining one inverse.  Therefore \(B\) is reduced and
finitely generated over \(K\).

Pull \(y_0\) back to \(B\) by TDR3, then change its bases by
the universal matrices and its section by \(s\), using TDR4.
This defines \(u\), with the source \texttt{UnivFrame} witness.
For (TDR.8), first pull \(y_0\) back by \(\psi\).  TDR4
gives the unique matrices and scalar comparing this reference
frame with \(z\).  Use precisely those coordinates in the
universal property of (TDR.7).  TDR3 and the compatibility in
TDR4 show that the resulting map carries \(u\) to \(z\).
If another map does so, uniqueness of the comparison coordinates
forces its values on every variable and on \(R\); it is therefore
the same map.  If \(m_\tau=0\), its determinant is one and it
contributes no variables, as required by the source.
\end{proof}
% END TDR6

% BEGIN TDR7
\paragraph{TDR7. Representability uses the actual theta and Lie frame bundles.}
\label{claim:TDR7}
The original functor \(S\mapsto\operatorname{FrPt}(S)\), with
\texttt{FrMap} as base change, is represented by
\[
 \mathcal Y=
 \left(\prod_{\tau,\,X}
   \operatorname{Isom}_X(\mathcal O_X^{m_\tau},V_\tau)\right)
       \times_X\operatorname{Isom}_X(\mathcal O_X,\Theta).
                                                               \tag{TDR.9}
\]
The morphism \(\mathcal Y\to X\) is affine and of finite
presentation.  In particular \(\mathcal Y\) is separated and
locally of finite type over \(K\).

\begin{proof}
On an open where a vector bundle of rank \(m\) is trivial,
its frame scheme is \(\operatorname{Spec}R[X_{ij},1/\det X]\).
On overlaps the invertible transition matrix changes these
coordinates by matrix multiplication, an algebraic isomorphism.
The cocycle law for the bundle is the gluing law for these schemes.
This constructs the indicated \(\operatorname{Isom}\) schemes.
The line factor is constructed in exactly the same way using
\(R[s,s^{-1}]\).  All these morphisms are affine and of finite
presentation, properties which can be checked on an open cover
of their base.  Their finite fiber product has the same properties;
on the common trivializing charts its algebra is precisely (TDR.7).

An \(S\)-point of (TDR.9) consists of a map \(x\) and the
listed trivializations on \(\operatorname{Spec}S\).  On an
affine base these are exactly the global ordered bases and
section generator in TDR1--TDR2.  Thus the correspondence with
the literal source record is bijective.  Its compatibility with
algebra maps is the actual \texttt{FrMap} relation, by TDR3;
no choice of a scalar or of a new model remains in this comparison.
Finally \(X\) is quasiprojective over \(K\), hence separated,
and \(\mathcal Y\to X\) is affine, hence separated.  Their
composition is separated.  Smoothness makes \(X\) locally of
finite presentation over \(K\), and the constructed morphism
has finite presentation, proving the final local finite-type
assertion.  This representability holds even for nonreduced
test coefficient algebras; the reduced-scheme representing
axiom of \texttt{Moduli} has not been extended or used outside
its original scope.
\end{proof}
% END TDR7

\subsection{Finite coefficients, without flatness assumptions}

% BEGIN TDR8
\paragraph{TDR8. Localization embeddings and simultaneous coefficient capture.}
\label{claim:TDR8}
Let \(S\) be any commutative \(K\)-algebra,
\(S_0\subseteq S\) a finitely generated subalgebra, and
\(h_1,\ldots,h_n\in S\) generate the unit ideal.
For each \(i\), let \(B_i\) be a finitely generated
\(K\)-algebra and \(a_i:B_i\to S_{h_i}\) an algebra map.
There is a finitely generated subalgebra \(S_1\subseteq S\)
containing \(S_0\) and all \(h_i\), in which the \(h_i\)
still generate the unit ideal, and maps
\[
 a_i':B_i\longrightarrow (S_1)_{h_i},\qquad
 \kappa_i:(S_1)_{h_i}\lhook\joinrel\longrightarrow S_{h_i},
 \qquad \kappa_i a_i'=a_i.                                  \tag{TDR.10}
\]
The \(\kappa_i\) are the canonical injective maps, with no
assumption that \(S_1\to S\) is flat or that the \(h_i\)
are non-zero-divisors.

\begin{proof}
First, if \(R\subseteq S\) and \(h\in R\), the canonical
map \(R_h\to S_h\) is injective.  If \(r/h^m\) maps to
zero, then \(h^N r=0\) in \(S\) for some \(N\).
The same equality holds in the subring \(R\), so \(r/h^m=0\)
in \(R_h\).  This argument includes nilpotents and zero divisors.

Choose \(c_i\in S\) with \(\sum_i c_i h_i=1\).
For each \(i\), choose a finite algebra generating set \(G_i\)
of \(B_i\).  Write
\(a_i(b)=s_{i,b}/h_i^{N_{i,b}}\) for \(b\in G_i\).
Adjoin to \(S_0\) the finitely many \(h_i,c_i,s_{i,b}\)
and call the resulting subalgebra \(S_1\).  The very same
equation \(\sum_i c_i h_i=1\) holds there, so its principal
opens still cover.  Every chosen generator image belongs to the
image of \(\kappa_i\), which is a subalgebra.  Hence the entire
image of \(a_i\) lies in that image.  Invert \(\kappa_i\)
on its image to obtain \(a_i'\); this is an algebra homomorphism
because \(\kappa_i\) is an injective algebra homomorphism.

Equivalently one may send the generators to the displayed fractions.
All relations automatically hold: their values map to zero under
the injective \(\kappa_i\), hence are already zero.  This is why
the argument needs only finite generation of the \(B_i\), not
an additional finite list of relations.  For an empty cover the
unit-ideal equation forces \(S\) to be the zero ring; the unique
zero subalgebra and empty family satisfy the assertion.
\end{proof}
% END TDR8

% BEGIN TDR9
\paragraph{TDR9. Principal-cover gluing of the original frames.}
\label{claim:TDR9}
Fix \(x:\operatorname{Spec}S\to X\) and a finite principal
cover \(D(h_i)\) with \((h_i)=S\).  Frames over its members,
with underlying points the restrictions of \(x\), glue uniquely
when their \texttt{FrMap} pullbacks agree over
\(S_{h_i h_j}\).  Equality of two frames over \(x\) can also
be checked after restriction to these localizations.

\begin{proof}
For every \(\tau,i\), the local tangent vectors are sections
of \(x^*V_\tau\).  Compatibility glues them by the sheaf axiom.
Their local inverse coordinate maps glue as well, or, equivalently,
the resulting map \(\mathcal O^{m_\tau}\to x^*V_\tau\)
is an isomorphism because it is one on each member of the cover.
This gives exactly the required ordered eigenbases by TDR1.
The local theta sections glue as sections of \(x^*\Theta\),
and their local inverses show that the glued section is a frame
of that line.  By TDR2 its section on \(A_x\) satisfies the
source \texttt{gen} condition.  All pullback equations are those
of TDR3, so the result is a frame of the original functor.

For clarity, the coefficient uniqueness implicit here works over
arbitrary \(S\).  If \(b\in S\) becomes zero in every
\(S_{h_i}\), choose \(N_i\) with \(h_i^{N_i}b=0\).
The ideal generated by these powers has radical the unit ideal,
and hence is itself the unit ideal.  A linear combination of the
powers is one, so \(b=0\).  Thus local equality of all matrix
entries and of the section scalar in TDR4 implies global equality.
Existence of glued module sections is the ordinary affine sheaf
equalizer, which follows by clearing finitely many denominators
and the same unit-ideal argument.  The proof does not require
faithful flatness of any individual localization map.
\end{proof}
% END TDR9

% BEGIN TDR10
\paragraph{TDR10. Finite-subalgebra descent of an actual scheme point.}
\label{claim:TDR10}
Let \(Z\) be a separated, locally finite-type \(K\)-scheme.
Every map \(x:\operatorname{Spec}S\to Z\), for an arbitrary
commutative \(K\)-algebra \(S\), is the base change of a map
\(x_0:\operatorname{Spec}S_0\to Z\) for a finitely generated
\(K\)-subalgebra \(S_0\subseteq S\).  A prescribed finite
set of elements of \(S\) may be included in \(S_0\).

\begin{proof}
Choose affine open neighborhoods \(U_i=\operatorname{Spec}B_i\)
of the image points, with \(B_i\) finitely generated over \(K\).
The inverse images cover \(\operatorname{Spec}S\), and its
principal opens form a basis.  Quasicompactness therefore gives
a finite principal cover \(D(h_i)\) such that
\(x(D(h_i))\subseteq U_i\).  We have maps
\(a_i:B_i\to S_{h_i}\).  Choose coefficients of an equation
\(\sum_i c_i h_i=1\), and use TDR8 to capture all the maps
in localizations of a single finitely generated subalgebra \(S_0\).
Start that construction with the algebra generated by the
prescribed finite set.  It gives maps
\(x_i:D_{S_0}(h_i)\to U_i\to Z\).

They agree on overlaps.  To see this without assuming flatness,
put \(R=(S_0)_{h_i h_j}\) and \(T=S_{h_i h_j}\).
The two overlap maps to \(Z\) coincide after base change to
\(T\), since both give the original \(x\).
As \(Z\) is separated over \(K\), their equalizer is the
pullback of the closed diagonal, hence is a closed subscheme of
\(\operatorname{Spec}R\), say defined by an ideal \(J\subseteq R\).
After base change to \(T\) the ideal is zero; thus every element
of \(J\) maps to zero in \(T\).  TDR8 says \(R\to T\)
is injective, so \(J=0\).  The two maps consequently agree
already over \(R\).  The maps \(x_i\) glue to \(x_0\),
and their base changes glue to \(x\).

The zero ring is handled by its empty spectrum.  This proof uses
only finite data, and in particular it does not impose a
Noetherian or reducedness assumption on \(S\).  It applies to
\(Z=X\) because the supplied moduli scheme is quasiprojective,
hence separated, and smooth, hence locally finite type.
It also applies directly to the actual frame scheme \(\mathcal Y\)
of TDR7.  The source's more general locally-finite-presentation
factorization theorem in \texttt{ModuliPointFromFgSubalgebra}
is therefore available here at every instance actually used.
\end{proof}
% END TDR10

% BEGIN TDR11
\paragraph{TDR11. Full finitely generated descent of every original framed point.}
\label{claim:TDR11}
Assuming exactly \(\mathrm{BC}\) and the supplied original data,
(TDR.1) holds.  This is the complete original
\texttt{TestDescent.exists\_fg\_descent} theorem.

\begin{proof}
There is a short proof: regard \(y\) as an \(S\)-point of
the explicitly constructed \(\mathcal Y\), use TDR10 for that
separated locally finite-type scheme, and apply the exact functor
identification TDR7.  It gives the actual frame \(y_1\), and
its base-change equality is precisely \texttt{FrMap} by TDR3.
We give the source-compatible local construction as well, to
record explicitly which finite data are retained and how the
underlying point and frames are joined.

First descend the underlying map of \(y\) to
\(x_0:\operatorname{Spec}S_0\to X\) by TDR10.
Choose a finite principal cover \(D(g_i)\) of
\(\operatorname{Spec}S_0\) such that each restriction of
\(x_0\) factors through a chart \(U_i=\operatorname{Spec}R_i\)
of TDR5.  Let \(h_i\) be the same elements in \(S\).
They still generate its unit ideal.  Let \(y_i\) be the
pullback of \(y\) to \(S_{h_i}\), and let
\(\psi_i:R_i\to S_{h_i}\) be the map supplied by the chart
factorization.  TDR6 supplies a universal frame \(u_i\) over
an explicit algebra \(B_i\) and a unique map
\[
 a_i:B_i\to S_{h_i},\qquad
 a_i|_{R_i}=\psi_i,\qquad(a_i)_*u_i=y_i.                   \tag{TDR.11}
\]

Apply TDR8 simultaneously to these finitely many \(a_i\),
starting with \(S_0\).  It gives one finitely generated
\(S_1\subseteq S\) and maps \(a_i'\) as in (TDR.10).
Let \(v_i=(a_i')_*u_i\) over \((S_1)_{h_i}\).
Its underlying point is exactly the restriction of
\(x_1=x_0\circ\operatorname{Spec}(S_1\supseteq S_0)\).
Indeed the chart map
\(R_i\to(S_0)_{g_i}\to(S_1)_{h_i}\) and the restriction
of \(a_i'\) to \(R_i\) become \(\psi_i\) in \(S_{h_i}\)
by (TDR.11).  The canonical localization map is injective,
so these two algebra maps already agree.  This proves equality
of the actual underlying maps, before attempting to glue frames.

On a double overlap \((S_1)_{h_i h_j}\), the two restrictions
of \(v_i,v_j\) have that same underlying map and become the
restriction of \(y\) after the canonical injection into
\(S_{h_i h_j}\).  TDR4 therefore makes them equal on the
original overlap.  TDR9 glues them to
\(y_1\in\operatorname{FrPt}(S_1)\).  Its base change to
\(S\) and \(y\) agree on the principal cover \(D(h_i)\),
so TDR9 gives their global equality.  In the source notation this
is the literal witness
\(\operatorname{FrMap}(S_1.\mathrm{val},y_1,y)\).

Every choice was made from the given \(y\).  There is no
independent choice of a new moduli object, no scalar rescaling
after the section is descended, and no assertion that the
inclusion \(S_1\subseteq S\) is flat.  Nilpotents in \(S\)
are retained if they occur among the finite coefficients.
\end{proof}
% END TDR11

% BEGIN TDR12
\paragraph{TDR12. Reduced finite-type algebras are exactly the original test-ring presentations.}
\label{claim:TDR12}
A commutative \(K\)-algebra \(R\) is reduced and finitely
generated if and only if it is isomorphic to
\(p.R=K[t_1,\ldots,t_n]/I\) for some literal
\(p:\texttt{TRing}\ K\), where \(I\) is radical.

\begin{proof}
Choose finitely many algebra generators to obtain a surjective
\(K\)-algebra map \(f:K[t_1,\ldots,t_n]\to R\).
If \(R\) is reduced and \(u^N\in\ker f\), then
\(f(u)^N=0\), whence \(f(u)=0\).  Thus \(I=\ker f\)
is radical, and the first isomorphism theorem gives the specified
algebra isomorphism \(p.R\simeq R\).  Conversely the images
of the variables generate the quotient; if a quotient class has
a nilpotent power, the radical condition makes that class zero.
This proof includes the unit ideal and the zero algebra, and
does not require a field-valued test point or a prime ideal.
\end{proof}
% END TDR12

% BEGIN TDR13
\paragraph{TDR13. The original AllTest theorem and the complex-frame return.}
\label{claim:TDR13}
Under \(\mathrm{BC}\), for every original \(E,H,D,l,X_o,c\),
the original \(\texttt{AllTest}\ X_o\ c\) holds.  In
particular, for every embedding \(\sigma:K\to\mathbb C\)
and every actual \(y\in\operatorname{FrPt}(\mathbb C_\sigma)\),
the exact source \texttt{FromTest} predicate holds.
The assertion applies simultaneously to every member of any
already chosen family of complex framed points.

\begin{proof}
Let \(S\) be reduced and let \(y\in\operatorname{FrPt}(S)\).
TDR11 gives a finitely generated subalgebra \(S_1\subseteq S\),
an actual frame \(y_1\), and
\(\operatorname{FrMap}(\iota,y_1,y)\), where \(\iota\)
is inclusion.  The subalgebra is reduced: a nilpotent element
would map to a nilpotent in \(S\), hence to zero, and inclusion
is injective.  TDR12 gives \(p:\texttt{TRing}\ K\) and
an algebra isomorphism \(e:p.R\xrightarrow{\sim}S_1\).
Pull \(y_1\) back by \(e^{-1}\) to obtain the literal
\(y'\in\operatorname{FrPt}(p.R)\), using TDR3.  Define
\(\phi=\iota\circ e\).  Functoriality gives
\[
 \phi_*y'=\iota_*e_*(e^{-1})_*y_1
         =\iota_*y_1=y.                                    \tag{TDR.12}
\]
Thus \((p,\phi,y')\) is exactly the triple required by
\texttt{FromTest}, proving (TDR.2) and the full universally
quantified \texttt{AllTest} statement.

The coefficient algebra \(\mathbb C_\sigma\), denoted
\texttt{CAt}\(\,\sigma\) in the source, is the ordinary
reduced ring \(\mathbb C\) with its \(K\)-structure given
by \(\sigma\).  Apply the result to it.  This does not alter
the embedding, the underlying point, its chosen eigenbases or
its chosen section.  Since the theorem holds for every such
frame, it applies to every \(\operatorname{un}(k,x)\) in
an already chosen family.  The test-ring witness may depend
on \((k,x)\); this is the original pointwise quantifier in
\texttt{exists\_unf\_hasPrinc}, not a demand for one test
ring presenting all members at once.
\end{proof}
% END TDR13

\subsection{Exact interfaces and the remaining geometry}

The endpoint \texttt{TestDescent.exists\_fg\_descent} is supplied
by TDR11, for every commutative coefficient algebra, without
reducedness.  The endpoint
\texttt{Frames.allTest\_of\_secBaseChange} is supplied by TDR13,
for every reduced coefficient algebra.  The source's intermediate
open-chart, universal-frame, frame-classification, injective
base-change equality and principal-cover gluing results are
proved in TDR3--TDR9, rather than left as new supplier assumptions.
TDR7 records why the same concrete frame functor is used throughout.

The hypothesis \(\mathrm{BC}\) remains named.  This chapter does
not prove its original instance from the full package of abelian
scheme laws, and does not assert existence of the supplied
\texttt{Moduli} record.  No \texttt{SecTwistSt} premise was needed.
The remaining normalized analytic frame family still has to
satisfy the exact \texttt{FrameOK}, Borel-weight, arithmetic
invariance and holomorphic-reading conditions.  TDR13 supplies its
\texttt{FromTest} field as soon as that same family has been
constructed.  The covering/fiber theorems, \texttt{IsUnf},
\texttt{exists\_unf\_hasPrinc}, spanning and common-family
obligations are not inferred from test-ring descent alone.

\paragraph{Pinned source correspondence.}
The primary sources are the pinned algebraic definitions of
\texttt{FrPt}, \texttt{FrMap}, \texttt{TRing} and
\texttt{AllTest}, the unitary-existence proof files, and the
original section-base-change predicate.  The accompanying
\texttt{source\_manifest.json} supplies complete immutable paths,
exact line ranges and file hashes; the claim provenance records
map every TDR claim to its source declarations.

The principal endpoint is in
\path{src/Development/AbelianVarietyTheory/Proofs/UnitaryExistence/FramedPointReducedFromTestRing.lean},
at lines 227--269 for finite-subalgebra descent and 276--290
for the return to all test rings.  TDR1--TDR7 prove the geometric
and universal-frame suppliers; TDR8--TDR11 prove finite coefficient
capture and gluing; TDR12--TDR13 perform the radical-quotient
presentation and exact source return.  The optional representability
shortcut in TDR7 and TDR11 is a human simplification; its literal
identification with the source frame records is proved before use.
